\chapter{应用数学：概率论}

概率论研究随机现象中的数量规律。它既可以从古典组合问题出发，也可以在测度论框架中被严格公理化。前者适合处理有限样本空间和计数问题，后者则能统一描述离散、连续以及更一般的随机对象。

本章先介绍经典概率、条件概率和随机变量，再过渡到测度论基础、条件期望、概率不等式、生成函数、收敛方式、极限定理以及随机过程的入门概念。

\section{经典概率与计数}

\subsection{随机试验与样本空间}

\begin{definition}[随机试验]
在相同条件下可以重复进行，并且结果具有不确定性的过程称为\textbf{随机试验}。
\end{definition}

\begin{definition}[样本空间与事件]
随机试验所有可能结果构成的集合称为\textbf{样本空间}，通常记作 $\Omega$。样本空间的子集称为\textbf{事件}。
\end{definition}

\begin{example}
投掷一枚硬币一次，样本空间可写为
\[
    \Omega=\{H,T\}.
\]
投掷一枚骰子一次，样本空间可写为
\[
    \Omega=\{1,2,3,4,5,6\}.
\]
事件“点数为偶数”就是
\[
    A=\{2,4,6\}.
\]
\end{example}

\subsection{概率的古典定义}

\begin{definition}[古典概率]
若样本空间 $\Omega$ 有限，并且所有基本结果等可能，则事件 $A\subseteq\Omega$ 的概率定义为
\[
    P(A)=\frac{|A|}{|\Omega|}.
\]
\end{definition}

\begin{example}
从一副标准扑克牌中随机抽一张，抽到红桃的概率为
\[
    \frac{13}{52}=\frac14.
\]
\end{example}

\subsection{基本计数原理}

\begin{proposition}[加法原理]
若两个任务不能同时发生，且第一个任务有 $m$ 种方法，第二个任务有 $n$ 种方法，则完成其中一个任务共有 $m+n$ 种方法。
\end{proposition}

\begin{proposition}[乘法原理]
若一个过程分为两个阶段，第一阶段有 $m$ 种选择，第二阶段对每种第一阶段选择都有 $n$ 种选择，则全过程共有 $mn$ 种选择。
\end{proposition}

\begin{definition}[排列]
从 $n$ 个不同对象中取出 $r$ 个并排列，方法数为
\[
    P(n,r)=\frac{n!}{(n-r)!}.
\]
\end{definition}

\begin{definition}[组合]
从 $n$ 个不同对象中取出 $r$ 个但不考虑顺序，方法数为
\[
    \binom nr=\frac{n!}{r!(n-r)!}.
\]
\end{definition}

\begin{example}
从 $10$ 个人中选出 $3$ 人组成委员会，方法数为
\[
    \binom{10}{3}=120.
\]
\end{example}

\begin{theorem}[二项式定理]
对任意非负整数 $n$，有
\[
    (x+y)^n=\sum_{k=0}^n\binom nk x^{n-k}y^k.
\]
\end{theorem}

\begin{remark}
组合数不仅用于计数，也自然出现在二项分布、生成函数和多项式展开中。
\end{remark}

\section{条件概率与独立性}

\subsection{条件概率}

\begin{definition}[条件概率]
若 $P(B)>0$，则在事件 $B$ 已发生的条件下，事件 $A$ 的条件概率定义为
\[
    P(A\mid B)=\frac{P(A\cap B)}{P(B)}.
\]
\end{definition}

\begin{example}
从一副扑克牌中抽一张。令 $A$ 为“抽到 A”，$B$ 为“抽到黑桃”。则
\[
    P(A\mid B)=\frac{1}{13}.
\]
\end{example}

\begin{proposition}[乘法公式]
若 $P(B)>0$，则
\[
    P(A\cap B)=P(A\mid B)P(B).
\]
更一般地，
\[
    P(A_1\cap\cdots\cap A_n)
    =P(A_1)P(A_2\mid A_1)\cdots P(A_n\mid A_1\cap\cdots\cap A_{n-1}).
\]
\end{proposition}

\subsection{独立性}

\begin{definition}[事件独立]
事件 $A$ 与 $B$ 称为独立，如果
\[
    P(A\cap B)=P(A)P(B).
\]
当 $P(B)>0$ 时，这等价于 $P(A\mid B)=P(A)$。
\end{definition}

\begin{definition}[相互独立]
事件族 $A_1,\dots,A_n$ 称为相互独立，如果对任意非空指标集 $I\subseteq\{1,\dots,n\}$，都有
\[
    P\left(\bigcap_{i\in I}A_i\right)=\prod_{i\in I}P(A_i).
\]
\end{definition}

\begin{remark}
两两独立不一定推出相互独立。相互独立要求任意多个事件的交都满足乘法关系。
\end{remark}

\subsection{Bayes 定理}

\begin{theorem}[全概率公式]
设 $B_1,\dots,B_n$ 构成样本空间的一个划分，且 $P(B_i)>0$。则对任意事件 $A$，有
\[
    P(A)=\sum_{i=1}^n P(A\mid B_i)P(B_i).
\]
\end{theorem}

\begin{theorem}[Bayes 定理]
在全概率公式的条件下，若 $P(A)>0$，则
\[
    P(B_j\mid A)=
    \frac{P(A\mid B_j)P(B_j)}{\sum_{i=1}^n P(A\mid B_i)P(B_i)}.
\]
\end{theorem}

\begin{remark}
Bayes 定理说明，观察到证据 $A$ 之后，先验概率 $P(B_j)$ 可以被更新为后验概率 $P(B_j\mid A)$。它是统计推断、机器学习和信息处理中的基础公式。
\end{remark}

\section{离散随机变量}

\subsection{定义}

\begin{definition}[离散随机变量]
随机变量是从样本空间到实数集的函数
\[
    X:\Omega\to\mathbb{R}.
\]
若 $X$ 的取值集合有限或可数，则称 $X$ 为离散随机变量。
\end{definition}

\begin{definition}[概率质量函数]
离散随机变量 $X$ 的概率质量函数定义为
\[
    p_X(x)=P(X=x).
\]
它满足 $p_X(x)\ge0$ 且
\[
    \sum_x p_X(x)=1.
\]
\end{definition}

\begin{definition}[分布函数]
随机变量 $X$ 的分布函数定义为
\[
    F_X(x)=P(X\le x).
\]
\end{definition}

\subsection{离散情形中的期望与方差}

\begin{definition}[期望]
离散随机变量 $X$ 的期望定义为
\[
    \mathbb{E}X=\sum_x xP(X=x),
\]
只要该级数绝对收敛。
\end{definition}

\begin{definition}[方差]
随机变量 $X$ 的方差定义为
\[
    \operatorname{Var}(X)=\mathbb{E}\left[(X-\mathbb{E}X)^2\right].
\]
\end{definition}

\begin{proposition}[方差的计算公式]
若 $\mathbb{E}X^2<\infty$，则
\[
    \operatorname{Var}(X)=\mathbb{E}X^2-(\mathbb{E}X)^2.
\]
\end{proposition}

\begin{proposition}[期望的线性性]
若期望存在，则
\[
    \mathbb{E}(aX+bY)=a\mathbb{E}X+b\mathbb{E}Y.
\]
这里不要求 $X,Y$ 独立。
\end{proposition}

\begin{remark}
期望线性性是概率论中最常用的工具之一。许多复杂随机变量的期望可以通过拆成指标变量来计算。
\end{remark}

\subsection{常见离散分布}

\begin{definition}[Bernoulli 分布]
若随机变量 $X$ 只取 $0$ 和 $1$，且
\[
    P(X=1)=p,
    \qquad
    P(X=0)=1-p,
\]
则称 $X\sim\operatorname{Bernoulli}(p)$。此时
\[
    \mathbb{E}X=p,
    \qquad
    \operatorname{Var}(X)=p(1-p).
\]
\end{definition}

\begin{definition}[二项分布]
若 $X$ 表示 $n$ 次独立 Bernoulli$(p)$ 试验中成功次数，则
\[
    X\sim\operatorname{Bin}(n,p),
\]
并且
\[
    P(X=k)=\binom nk p^k(1-p)^{n-k},
    \qquad k=0,1,\dots,n.
\]
此时
\[
    \mathbb{E}X=np,
    \qquad
    \operatorname{Var}(X)=np(1-p).
\]
\end{definition}

\begin{example}
投掷一枚不均匀硬币 $n$ 次，正面出现次数服从二项分布。
\end{example}

\begin{definition}[几何分布]
若 $X$ 表示第一次成功出现所需试验次数，且每次成功概率为 $p$，则
\[
    P(X=k)=(1-p)^{k-1}p,
    \qquad k=1,2,\dots.
\]
称 $X$ 服从几何分布。
\end{definition}

\begin{definition}[Poisson 分布]
若
\[
    P(X=k)=e^{-\lambda}\frac{\lambda^k}{k!},
    \qquad k=0,1,2,\dots,
\]
其中 $\lambda>0$，则称 $X\sim\operatorname{Poisson}(\lambda)$。此时
\[
    \mathbb{E}X=\lambda,
    \qquad
    \operatorname{Var}(X)=\lambda.
\]
\end{definition}

\begin{remark}
Poisson 分布常用于近似单位时间或单位区域内稀有事件出现的次数。
\end{remark}

\section{连续随机变量}

\subsection{密度函数}

\begin{definition}[概率密度函数]
随机变量 $X$ 称为连续型，如果存在非负函数 $f_X$，使得对任意集合 $A$，有
\[
    P(X\in A)=\int_A f_X(x)\,dx.
\]
函数 $f_X$ 称为 $X$ 的概率密度函数，并满足
\[
    \int_{-\infty}^{\infty}f_X(x)\,dx=1.
\]
\end{definition}

\begin{definition}[连续变量的分布函数]
若 $X$ 有密度 $f_X$，则
\[
    F_X(x)=P(X\le x)=\int_{-\infty}^x f_X(t)\,dt.
\]
在连续点处，$F_X'(x)=f_X(x)$。
\end{definition}

\begin{definition}[期望]
连续随机变量的期望定义为
\[
    \mathbb{E}X=\int_{-\infty}^{\infty}x f_X(x)\,dx,
\]
只要积分绝对收敛。
\end{definition}

\subsection{常见连续分布}

\begin{definition}[均匀分布]
若 $X$ 在区间 $[a,b]$ 上具有密度
\[
    f_X(x)=
    \begin{cases}
    \dfrac1{b-a},&a\le x\le b,\\
    0,&\text{否则},
    \end{cases}
\]
则称 $X\sim U(a,b)$。此时
\[
    \mathbb{E}X=\frac{a+b}{2},
    \qquad
    \operatorname{Var}(X)=\frac{(b-a)^2}{12}.
\]
\end{definition}

\begin{definition}[指数分布]
若 $X$ 的密度为
\[
    f_X(x)=
    \begin{cases}
    \lambda e^{-\lambda x},&x\ge0,\\
    0,&x<0,
    \end{cases}
\]
其中 $\lambda>0$，则称 $X\sim\operatorname{Exp}(\lambda)$。
\end{definition}

\begin{definition}[正态分布]
若 $X$ 的密度为
\[
    f_X(x)=\frac1{\sqrt{2\pi}\sigma}
    \exp\left(-\frac{(x-\mu)^2}{2\sigma^2}\right),
\]
其中 $\sigma>0$，则称 $X\sim N(\mu,\sigma^2)$。
\end{definition}

\begin{remark}
正态分布之所以重要，一个根本原因是中心极限定理：大量独立小扰动的和，在适当标准化后往往趋向正态分布。
\end{remark}

\section{测度论基础}

\subsection{为什么经典概率还不够}

古典概率依赖有限且等可能的样本空间；连续概率依赖积分；随机过程则需要在函数空间上定义概率。为了把这些情况放在统一框架中，需要使用 $\sigma$-代数和概率测度。

\subsection{$\sigma$-代数}

\begin{definition}[$\sigma$-代数]
设 $\Omega$ 是集合。子集族 $\mathcal{F}\subseteq\mathcal{P}(\Omega)$ 称为 $\sigma$-代数，如果：
\begin{enumerate}
    \item $\Omega\in\mathcal{F}$；
    \item 若 $A\in\mathcal{F}$，则 $A^c\in\mathcal{F}$；
    \item 若 $A_1,A_2,\dots\in\mathcal{F}$，则
    \[
        \bigcup_{n=1}^{\infty}A_n\in\mathcal{F}.
    \]
\end{enumerate}
$\mathcal{F}$ 中的集合称为可测事件。
\end{definition}

\begin{proposition}
若 $\mathcal{F}$ 是 $\sigma$-代数，则它对可数交也封闭，即
\[
    \bigcap_{n=1}^{\infty}A_n\in\mathcal{F}
\]
只要每个 $A_n\in\mathcal{F}$。
\end{proposition}

\begin{definition}[Borel $\sigma$-代数]
拓扑空间 $X$ 中由所有开集生成的最小 $\sigma$-代数称为 Borel $\sigma$-代数，记作 $\mathcal{B}(X)$。
\end{definition}

\subsection{概率空间}

\begin{definition}[概率测度]
设 $\mathcal{F}$ 是 $\Omega$ 上的 $\sigma$-代数。函数 $P:\mathcal{F}\to[0,1]$ 称为概率测度，如果
\[
    P(\Omega)=1
\]
且对两两不交的事件 $A_1,A_2,\dots$，有可列可加性
\[
    P\left(\bigcup_{n=1}^{\infty}A_n\right)=\sum_{n=1}^{\infty}P(A_n).
\]
\end{definition}

\begin{definition}[概率空间]
三元组
\[
    (\Omega,\mathcal{F},P)
\]
称为\textbf{概率空间}。
\end{definition}

\begin{proposition}[基本性质]
对任意事件 $A,B\in\mathcal{F}$，有
\[
    P(A^c)=1-P(A),
\]
若 $A\subseteq B$，则 $P(A)\le P(B)$，并且
\[
    P(A\cup B)=P(A)+P(B)-P(A\cap B).
\]
\end{proposition}

\begin{proposition}[概率的连续性]
若 $A_1\subseteq A_2\subseteq\cdots$，则
\[
    P\left(\bigcup_{n=1}^{\infty}A_n\right)=\lim_{n\to\infty}P(A_n).
\]
若 $A_1\supseteq A_2\supseteq\cdots$，则
\[
    P\left(\bigcap_{n=1}^{\infty}A_n\right)=\lim_{n\to\infty}P(A_n).
\]
\end{proposition}

\section{随机变量作为可测函数}

\subsection{可测性}

\begin{definition}[随机变量]
设 $(\Omega,\mathcal{F},P)$ 是概率空间。函数 $X:\Omega\to\mathbb{R}$ 称为随机变量，如果对任意 Borel 集 $B\in\mathcal{B}(\mathbb{R})$，都有
\[
    X^{-1}(B)\in\mathcal{F}.
\]
也就是说，$X$ 是从 $(\Omega,\mathcal{F})$ 到 $(\mathbb{R},\mathcal{B}(\mathbb{R}))$ 的可测函数。
\end{definition}

\begin{remark}
可测性的意义在于，形如 $\{X\in B\}$ 的事件必须能被赋予概率。
\end{remark}

\begin{definition}[分布]
随机变量 $X$ 的分布是 $\mathbb{R}$ 上的概率测度 $\mu_X$，定义为
\[
    \mu_X(B)=P(X\in B)=P(X^{-1}(B)).
\]
\end{definition}

\begin{definition}[随机向量]
映射 $X:\Omega\to\mathbb{R}^n$ 若可测，则称为随机向量。它的分布是 $\mathbb{R}^n$ 上的概率测度。
\end{definition}

\subsection{作为积分的期望}

\begin{definition}[期望]
随机变量 $X$ 的期望定义为 Lebesgue 积分
\[
    \mathbb{E}X=\int_\Omega X\,dP,
\]
只要该积分有意义。
\end{definition}

\begin{theorem}[无意识统计学家定律]
若 $X$ 的分布为 $\mu_X$，而 $g$ 是可测函数，则
\[
    \mathbb{E}g(X)=\int_{\mathbb{R}} g(x)\,d\mu_X(x),
\]
只要两边有意义。离散情形下为
\[
    \mathbb{E}g(X)=\sum_x g(x)P(X=x),
\]
连续密度情形下为
\[
    \mathbb{E}g(X)=\int_{-\infty}^{\infty}g(x)f_X(x)\,dx.
\]
\end{theorem}

\section{联合分布与独立性}

\subsection{联合分布}

\begin{definition}[联合分布]
随机向量 $(X,Y)$ 的分布称为 $X$ 与 $Y$ 的联合分布，定义为
\[
    \mu_{X,Y}(A)=P((X,Y)\in A),
    \qquad A\in\mathcal{B}(\mathbb{R}^2).
\]
\end{definition}

\begin{definition}[联合密度]
若存在非负函数 $f_{X,Y}$，使得
\[
    P((X,Y)\in A)=\iint_A f_{X,Y}(x,y)\,dxdy,
\]
则称 $f_{X,Y}$ 为联合密度。
\end{definition}

\begin{definition}[边缘密度]
若 $(X,Y)$ 有联合密度 $f_{X,Y}$，则边缘密度为
\[
    f_X(x)=\int_{-\infty}^{\infty}f_{X,Y}(x,y)\,dy,
\]
\[
    f_Y(y)=\int_{-\infty}^{\infty}f_{X,Y}(x,y)\,dx.
\]
\end{definition}

\subsection{随机变量的独立性}

\begin{definition}[随机变量独立]
随机变量 $X,Y$ 称为独立，如果对任意 Borel 集 $A,B$，有
\[
    P(X\in A,Y\in B)=P(X\in A)P(Y\in B).
\]
\end{definition}

\begin{proposition}[分解判别法]
若 $X,Y$ 有联合密度，则 $X,Y$ 独立，当且仅当
\[
    f_{X,Y}(x,y)=f_X(x)f_Y(y)
\]
几乎处处成立。离散情形中对应为
\[
    P(X=x,Y=y)=P(X=x)P(Y=y).
\]
\end{proposition}

\begin{proposition}
若 $X,Y$ 独立，且期望存在，则
\[
    \mathbb{E}(XY)=\mathbb{E}X\,\mathbb{E}Y.
\]
\end{proposition}

\begin{remark}
反过来，$\mathbb{E}(XY)=\mathbb{E}X\,\mathbb{E}Y$ 通常不能推出独立，只能说明二者不相关。
\end{remark}

\subsection{协方差与相关系数}

\begin{definition}[协方差]
随机变量 $X,Y$ 的协方差定义为
\[
    \operatorname{Cov}(X,Y)
    =\mathbb{E}\left[(X-\mathbb{E}X)(Y-\mathbb{E}Y)\right].
\]
\end{definition}

\begin{definition}[相关系数]
若 $\operatorname{Var}(X)>0$ 且 $\operatorname{Var}(Y)>0$，定义
\[
    \rho_{X,Y}=\frac{\operatorname{Cov}(X,Y)}{\sqrt{\operatorname{Var}(X)\operatorname{Var}(Y)}}.
\]
\end{definition}

\newpage

\begin{proposition}
相关系数满足
\[
    -1\le\rho_{X,Y}\le1.
\]
这来自 Cauchy--Schwarz 不等式。
\end{proposition}

\section{条件期望}

\subsection{关于事件的条件化}

若 $P(A)>0$，则
\[
    \mathbb{E}(X\mid A)=\frac{\mathbb{E}(X\mathbf{1}_A)}{P(A)}.
\]
它表示在事件 $A$ 已发生的条件下，对 $X$ 的平均值。

\subsection{关于随机变量的条件化}

对于离散随机变量，条件期望可以写作
\[
    \mathbb{E}(Y\mid X=x)=\sum_y yP(Y=y\mid X=x).
\]
于是 $\mathbb{E}(Y\mid X)$ 本身是 $X$ 的函数。

\begin{definition}[关于 $\sigma$-代数的条件期望]
设 $X$ 可积，$\mathcal{G}\subseteq\mathcal{F}$ 是子 $\sigma$-代数。随机变量 $Z$ 称为 $X$ 关于 $\mathcal{G}$ 的条件期望，记作
\[
    Z=\mathbb{E}(X\mid\mathcal{G}),
\]
如果 $Z$ 是 $\mathcal{G}$-可测的，并且对任意 $A\in\mathcal{G}$，都有
\[
    \int_A Z\,dP=\int_A X\,dP.
\]
\end{definition}

\begin{definition}[关于随机变量的条件期望]
定义
\[
    \mathbb{E}(Y\mid X)=\mathbb{E}(Y\mid\sigma(X)),
\]
其中 $\sigma(X)$ 是由 $X$ 生成的 $\sigma$-代数。
\end{definition}

\begin{theorem}[塔式性质]
若 $\mathcal{H}\subseteq\mathcal{G}\subseteq\mathcal{F}$，则
\[
    \mathbb{E}\left(\mathbb{E}(X\mid\mathcal{G})\mid\mathcal{H}\right)
    =\mathbb{E}(X\mid\mathcal{H}).
\]
特别地，
\[
    \mathbb{E}\left(\mathbb{E}(X\mid\mathcal{G})\right)=\mathbb{E}X.
\]
\end{theorem}

\begin{remark}
条件期望不是简单的数，而是给定信息后最好的可测预测。它在鞅论、统计估计和随机过程理论中十分核心。
\end{remark}

\section{概率不等式}

\begin{theorem}[Markov 不等式]
若 $X\ge0$ 且 $a>0$，则
\[
    P(X\ge a)\le\frac{\mathbb{E}X}{a}.
\]
\end{theorem}

\begin{theorem}[Chebyshev 不等式]
若 $X$ 有有限方差，则对任意 $\varepsilon>0$，有
\[
    P(|X-\mathbb{E}X|\ge\varepsilon)
    \le\frac{\operatorname{Var}(X)}{\varepsilon^2}.
\]
\end{theorem}

\begin{theorem}[Jensen 不等式]
若 $\varphi$ 是凸函数，且相关期望存在，则
\[
    \varphi(\mathbb{E}X)\le\mathbb{E}\varphi(X).
\]
\end{theorem}

\begin{theorem}[Cauchy--Schwarz 不等式]
若 $X,Y\in L^2$，则
\[
    |\mathbb{E}(XY)|^2\le\mathbb{E}(X^2)\mathbb{E}(Y^2).
\]
\end{theorem}

\section{生成函数与特征函数}

\subsection{概率生成函数}

\begin{definition}[概率生成函数]
若 $X$ 是取非负整数值的随机变量，则它的概率生成函数定义为
\[
    G_X(s)=\mathbb{E}s^X=\sum_{k=0}^{\infty}P(X=k)s^k.
\]
\end{definition}

\begin{proposition}
若 $X,Y$ 独立且取非负整数值，则
\[
    G_{X+Y}(s)=G_X(s)G_Y(s).
\]
\end{proposition}

\subsection{矩母函数}

\begin{definition}[矩母函数]
随机变量 $X$ 的矩母函数定义为
\[
    M_X(t)=\mathbb{E}e^{tX},
\]
只要该期望在 $t=0$ 附近存在。
\end{definition}

\begin{remark}
若矩母函数在零点附近存在，则其导数给出矩：
\[
    M_X^{(n)}(0)=\mathbb{E}X^n.
\]
\end{remark}

\subsection{特征函数}

\begin{definition}[特征函数]
随机变量 $X$ 的特征函数定义为
\[
    \varphi_X(t)=\mathbb{E}e^{itX},
    \qquad t\in\mathbb{R}.
\]
\end{definition}

\begin{proposition}[基本性质]
特征函数总是存在，并满足
\[
    \varphi_X(0)=1,
    \qquad
    |\varphi_X(t)|\le1.
\]
若 $X,Y$ 独立，则
\[
    \varphi_{X+Y}(t)=\varphi_X(t)\varphi_Y(t).
\]
\end{proposition}

\begin{theorem}[唯一性定理]
随机变量的特征函数唯一决定其分布。
\end{theorem}

\section{收敛方式}

\begin{definition}[几乎处处收敛]
若
\[
    P\left(\lim_{n\to\infty}X_n=X\right)=1,
\]
则称 $X_n$ 几乎处处收敛到 $X$，记作
\[
    X_n\to X\quad a.s.
\]
\end{definition}

\begin{definition}[依概率收敛]
若对任意 $\varepsilon>0$，有
\[
    P(|X_n-X|>\varepsilon)\to0,
\]
则称 $X_n$ 依概率收敛到 $X$，记作
\[
    X_n\xrightarrow{P}X.
\]
\end{definition}

\begin{definition}[$L^p$ 收敛]
若
\[
    \mathbb{E}|X_n-X|^p\to0,
\]
则称 $X_n$ 在 $L^p$ 中收敛到 $X$。
\end{definition}

\begin{definition}[依分布收敛]
若 $X_n$ 的分布函数 $F_n$ 在 $X$ 的分布函数 $F$ 的每个连续点处满足
\[
    F_n(x)\to F(x),
\]
则称 $X_n$ 依分布收敛到 $X$，记作
\[
    X_n\xrightarrow{d}X.
\]
\end{definition}

\begin{theorem}[收敛方式之间的关系]
有如下蕴含关系：
\[
    X_n\to X\ a.s.\quad\Longrightarrow\quad X_n\xrightarrow{P}X
    \quad\Longrightarrow\quad X_n\xrightarrow{d}X.
\]
若 $X_n\to X$ 于 $L^p$，则 $X_n\xrightarrow{P}X$。
\end{theorem}

\begin{remark}
依分布收敛是最弱的常见收敛方式之一，它只描述分布的极限，不要求随机变量在同一样本点上接近。
\end{remark}

\section{极限定理}

\subsection{大数定律}

\begin{theorem}[弱大数定律]
设 $X_1,X_2,\dots$ 独立同分布，且
\[
    \mathbb{E}X_i=\mu,
    \qquad
    \operatorname{Var}(X_i)=\sigma^2<\infty.
\]
令
\[
    \overline X_n=\frac1n\sum_{i=1}^n X_i.
\]
则
\[
    \overline X_n\xrightarrow{P}\mu.
\]
\end{theorem}

\begin{proof}
由期望线性性，$\mathbb{E}\overline X_n=\mu$；由独立性，
\[
    \operatorname{Var}(\overline X_n)=\frac{\sigma^2}{n}.
\]
对 $\overline X_n$ 使用 Chebyshev 不等式即可得到结论。
\end{proof}

\begin{theorem}[强大数定律]
若 $X_1,X_2,\dots$ 独立同分布且 $\mathbb{E}|X_1|<\infty$，则
\[
    \frac1n\sum_{i=1}^n X_i\to\mathbb{E}X_1
    \quad a.s.
\]
\end{theorem}

\begin{remark}
弱大数定律给出依概率收敛，强大数定律给出几乎处处收敛。后者更强，也更接近频率稳定性的直观表述。
\end{remark}

\subsection{中心极限定理}

\begin{theorem}[Lindeberg--Levy 中心极限定理]
设 $X_1,X_2,\dots$ 独立同分布，且
\[
    \mathbb{E}X_i=\mu,
    \qquad
    \operatorname{Var}(X_i)=\sigma^2>0.
\]
则
\[
    \frac{\sum_{i=1}^n X_i-n\mu}{\sigma\sqrt n}
    \xrightarrow{d}N(0,1).
\]
\end{theorem}

中心极限定理说明，独立同分布随机变量之和在标准化后趋向正态分布。这一结论解释了为什么正态分布在统计学和自然科学中频繁出现。

\begin{remark}
中心极限定理中的标准化有两个作用：减去均值使中心回到 $0$，除以标准差尺度 $\sigma\sqrt n$ 使方差稳定为 $1$。
\end{remark}

\section{Borel--Cantelli 引理}

\begin{definition}[事件的上极限]
事件列 $A_1,A_2,\dots$ 的上极限定义为
\[
    \limsup_{n\to\infty}A_n
    =\bigcap_{n=1}^{\infty}\bigcup_{k=n}^{\infty}A_k.
\]
它表示“无限多个 $A_n$ 发生”的事件。
\end{definition}

\newpage

\begin{theorem}[第一 Borel--Cantelli 引理]
若
\[
    \sum_{n=1}^{\infty}P(A_n)<\infty,
\]
则
\[
    P(A_n\ \text{无限多次发生})=0.
\]
即
\[
    P\left(\limsup_{n\to\infty}A_n\right)=0.
\]
\end{theorem}

\begin{theorem}[第二 Borel--Cantelli 引理]
若事件 $A_1,A_2,\dots$ 相互独立，并且
\[
    \sum_{n=1}^{\infty}P(A_n)=\infty,
\]
则
\[
    P(A_n\ \text{无限多次发生})=1.
\]
\end{theorem}

\section{随机过程导论}

\begin{definition}[随机过程]
随机过程是一族随机变量
\[
    \{X_t:t\in T\},
\]
其中 $T$ 称为指标集。若 $T=\mathbb{N}$，称为离散时间随机过程；若 $T=[0,\infty)$，称为连续时间随机过程。
\end{definition}

\begin{example}
简单随机游走、Poisson 过程、Brown 运动、马尔可夫链和时间序列模型都是随机过程的例子。
\end{example}

\begin{definition}[滤过]
一族 $\sigma$-代数 $\{\mathcal{F}_t\}_{t\in T}$ 称为滤过，如果当 $s\le t$ 时，
\[
    \mathcal{F}_s\subseteq\mathcal{F}_t.
\]
它表示随时间增长而增加的信息。
\end{definition}

\begin{definition}[鞅]
设 $\{X_t\}$ 适应滤过 $\{\mathcal{F}_t\}$，且 $\mathbb{E}|X_t|<\infty$。如果对 $s\le t$，有
\[
    \mathbb{E}(X_t\mid\mathcal{F}_s)=X_s,
\]
则称 $\{X_t\}$ 是鞅。
\end{definition}

\begin{remark}
鞅可以理解为“公平游戏”的数学模型：在已知当前信息的条件下，未来的期望等于现在的值。
\end{remark}

\begin{example}[简单随机游走]
设 $\xi_1,\xi_2,\dots$ 独立同分布，且
\[
    P(\xi_i=1)=P(\xi_i=-1)=\frac12.
\]
定义
\[
    S_n=\sum_{i=1}^n \xi_i.
\]
则 $\{S_n\}$ 是一维简单对称随机游走，并且相对于自然滤过是鞅。
\end{example}

\section{分布函数与随机变量变换}

概率分布可以用若干种等价方式描述。离散分布通常由概率质量函数给出，连续分布通常由密度函数给出，而一般的实值分布则可由累积分布函数给出。在这些描述中，分布函数最为普适，因为即使一个分布既非纯离散也非纯连续，它仍然存在。

\begin{definition}[累积分布函数]
对实值随机变量 $X$，其\textbf{累积分布函数}定义为
\[
    F_X(x)=\mathbb{P}(X\leq x).
\]
\end{definition}

\begin{theorem}[分布函数的刻画]
函数 $F:\mathbb{R}\to[0,1]$ 是某个实值随机变量的分布函数，当且仅当：
\begin{enumerate}
    \item $F$ 单调不减；
    \item $F$ 右连续；
    \item $\lim_{x\to-\infty}F(x)=0$ 且 $\lim_{x\to\infty}F(x)=1$。
\end{enumerate}
此外，
\[
    \mathbb{P}(X=x)=F_X(x)-F_X(x^-),
\]
所以分布函数的跳跃恰好对应点质量。
\end{theorem}

\begin{example}[混合分布]
设 $X$ 以概率 $1/2$ 取值 $0$，而以概率 $1/2$ 在 $[0,1]$ 上服从均匀分布。则
\[
F_X(x)=
\begin{cases}
0,&x<0,\\
\frac12+\frac{x}{2},&0\leq x<1,\\
1,&x\geq1.
\end{cases}
\]
在 $0$ 处的跳跃记录原子质量，线性部分则记录连续分量。不存在一个通常意义下的密度函数能够描述整个分布。
\end{example}

\subsection{单个随机变量的函数}

\begin{theorem}[单调变量变换]
设 $X$ 的密度为 $f_X$，令 $Y=g(X)$，其中 $g$ 严格单调且连续可微。则在 $g$ 的值域上，
\[
    f_Y(y)=f_X(g^{-1}(y))
    \left|\frac{d}{dy}g^{-1}(y)\right|.
\]
\end{theorem}

\begin{proof}
如果 $g$ 单调递增，则
\[
    F_Y(y)=\mathbb{P}(g(X)\leq y)
    =\mathbb{P}(X\leq g^{-1}(y))
    =F_X(g^{-1}(y)).
\]
对其求导并应用链式法则。若 $g$ 单调递减，不等号方向反转，最终得到同一个含绝对值的公式。
\end{proof}

\begin{example}[标准正态随机变量的平方]
设 $X\sim N(0,1)$，令 $Y=X^2$。映射 $x\mapsto x^2$ 不是一一映射，因此两个反函数分支都要贡献：
\[
    f_Y(y)
    =\frac{f_X(\sqrt y)}{2\sqrt y}
     +\frac{f_X(-\sqrt y)}{2\sqrt y}
    =\frac{1}{\sqrt{2\pi y}}e^{-y/2},
    \qquad y>0.
\]
因此 $X^2$ 服从自由度为 $1$ 的卡方分布。
\end{example}

\begin{theorem}[尾积分公式]
如果 $X\geq0$，则
\[
    \mathbb{E}[X]=\int_0^\infty\mathbb{P}(X>t)\,dt,
\]
其中等式两边都允许取无穷大。
\end{theorem}

\begin{proof}
逐点有
\[
    X(\omega)=\int_0^\infty
    \mathbf{1}_{\{X(\omega)>t\}}\,dt.
\]
由 Tonelli 定理可以交换积分次序：
\[
    \mathbb{E}[X]
    =\int_0^\infty
    \mathbb{E}[\mathbf{1}_{\{X>t\}}]dt
    =\int_0^\infty\mathbb{P}(X>t)dt.
\]
\end{proof}

\section{联合密度、卷积与条件化}

对多个随机变量而言，边缘分布分别描述各个坐标，却不能决定它们之间的依赖关系；联合分布才包含完整信息。边缘化通过积分消去不需要的坐标，条件化沿某个坐标重新归一化，而卷积描述独立随机变量之和的分布。

\begin{definition}[联合密度]
如果对每个适当的 Borel 集 $A\subseteq\mathbb{R}^2$，都有
\[
    \mathbb{P}((X,Y)\in A)
    =\iint_A f_{X,Y}(x,y)\,dx\,dy,
\]
则称随机向量 $(X,Y)$ 具有联合密度 $f_{X,Y}$。其边缘密度为
\[
    f_X(x)=\int_{-\infty}^{\infty}f_{X,Y}(x,y)\,dy,
    \qquad
    f_Y(y)=\int_{-\infty}^{\infty}f_{X,Y}(x,y)\,dx.
\]
\end{definition}

\begin{theorem}[独立性的密度判据]
如果 $(X,Y)$ 具有联合密度，那么 $X$ 与 $Y$ 独立，当且仅当
\[
    f_{X,Y}(x,y)=f_X(x)f_Y(y)
\]
对几乎处处的 $(x,y)$ 成立。
\end{theorem}

\begin{definition}[条件密度]
如果 $f_Y(y)>0$，定义
\[
    f_{X\mid Y}(x\mid y)
    =\frac{f_{X,Y}(x,y)}{f_Y(y)}.
\]
对固定的 $y$，这是关于 $x$ 的密度，表示在条件 $Y=y$ 下 $X$ 的分布。
\end{definition}

\begin{example}[三角形上的均匀分布]
设 $(X,Y)$ 在
\[
    T=\{(x,y):0<y<x<1\}
\]
上均匀分布。由于 $T$ 的面积为 $1/2$，联合密度在 $T$ 上等于 $2$。当 $0<x<1$ 时，
\[
    f_X(x)=\int_0^x2\,dy=2x.
\]
对固定的 $x$，给定 $X=x$ 时 $Y$ 的条件密度为
\[
    f_{Y\mid X}(y\mid x)=\frac{2}{2x}=\frac1x,
    \qquad0<y<x.
\]
因此，在条件 $X=x$ 下，$Y$ 在 $(0,x)$ 上均匀分布。
\end{example}

\begin{theorem}[卷积公式]
如果 $X,Y$ 是相互独立的连续随机变量，密度分别为 $f_X,f_Y$，那么 $S=X+Y$ 的密度为
\[
    f_S(s)=\int_{-\infty}^{\infty}f_X(x)f_Y(s-x)\,dx.
\]
\end{theorem}

\begin{proof}
变量变换 $(x,y)\mapsto(x,s=x+y)$ 的 Jacobian 行列式为 $1$。由于独立性给出联合密度 $f_X(x)f_Y(y)$，对 $x$ 积分消元即得到上述公式。
\end{proof}

\begin{example}[两个均匀随机变量之和]
设 $X,Y$ 相互独立，且都在 $[0,1]$ 上均匀分布。则
\[
    f_{X+Y}(s)
    =\int_{\max(0,s-1)}^{\min(1,s)}1\,dx
    =
    \begin{cases}
    s,&0<s<1,\\
    2-s,&1\leq s<2,\\
    0,&\text{其他情形}.
    \end{cases}
\]
这个三角形密度反映了直线段 $x+y=s$ 在单位正方形内部的长度先增大后减小。
\end{example}

\section{期望的收敛定理}

随机变量的极限只有在能够穿过概率或期望时才真正有用。测度论给出了这样做的精确条件。单调收敛定理处理递增的非负逼近，控制收敛定理则处理被某个可积函数统一控制的带符号随机变量。

\begin{theorem}[单调收敛定理]
如果 $0\leq X_1\leq X_2\leq\cdots$，且 $X_n\to X$ 几乎处处成立，那么
\[
    \mathbb{E}[X_n]\uparrow\mathbb{E}[X].
\]
\end{theorem}

\begin{theorem}[Fatou 引理]
如果 $X_n\geq0$，则
\[
    \mathbb{E}\left[\liminf_{n\to\infty}X_n\right]
    \leq
    \liminf_{n\to\infty}\mathbb{E}[X_n].
\]
\end{theorem}

\begin{theorem}[控制收敛定理]
设 $X_n\to X$ 几乎处处，并且存在可积随机变量 $Y$，使得对每个 $n$ 都几乎处处有
\[
    |X_n|\leq Y.
\]
则 $X$ 可积，并且
\[
    \mathbb{E}[X_n]\to\mathbb{E}[X].
\]
\end{theorem}

\begin{example}[为什么逐点收敛还不够]
在带 Lebesgue 概率测度的 $(0,1)$ 上，令
\[
    X_n=n\mathbf{1}_{(0,1/n)}.
\]
则 $X_n\to0$ 几乎处处，但对每个 $n$，
\[
    \mathbb{E}[X_n]=1.
\]
不存在一个可积随机变量同时控制所有 $X_n$。这个例子说明，要交换极限与期望，仅有逐点收敛并不够。
\end{example}

\begin{corollary}[特征函数的连续性]
对任意随机变量 $X$，其特征函数
\[
    \varphi_X(t)=\mathbb{E}[e^{itX}]
\]
关于 $t$ 连续。
\end{corollary}

\begin{proof}
如果 $t_n\to t$，则 $e^{it_nX}\to e^{itX}$ 几乎处处，同时
\[
    |e^{it_nX}|=1.
\]
以可积函数 $1$ 为控制函数，应用控制收敛定理即可。
\end{proof}

\section{作为信息的条件期望}

对事件进行条件化得到一个数；对 $\sigma$-代数进行条件化则得到一个随机变量，因为可用信息可能在不同状态下产生不同预测。条件期望 $\mathbb{E}[X\mid\mathcal{G}]$ 是平均意义下对 $X$ 的最佳 $\mathcal{G}$-可测概括。

\begin{definition}[条件期望]
设 $X$ 可积，$\mathcal{G}\subseteq\mathcal{F}$ 是子 $\sigma$-代数。如果随机变量 $Y$ 满足：
\begin{enumerate}
    \item $Y$ 是 $\mathcal{G}$-可测的；
    \item $Y$ 可积；
    \item 对每个 $G\in\mathcal{G}$，
    \[
        \int_GY\,d\mathbb{P}=\int_GX\,d\mathbb{P},
    \]
\end{enumerate}
则称 $Y$ 是 $\mathbb{E}[X\mid\mathcal{G}]$ 的一个版本。它在几乎处处相等的意义下唯一。
\end{definition}

\begin{theorem}[基本性质]
对可积随机变量 $X,Y$ 和常数 $a,b$：
\begin{enumerate}
    \item \textbf{线性性}：
    \[
        \mathbb{E}[aX+bY\mid\mathcal{G}]
        =a\mathbb{E}[X\mid\mathcal{G}]
        +b\mathbb{E}[Y\mid\mathcal{G}].
    \]
    \item \textbf{提出已知因子}：如果 $Z$ 有界且 $\mathcal{G}$-可测，则
    \[
        \mathbb{E}[ZX\mid\mathcal{G}]
        =Z\mathbb{E}[X\mid\mathcal{G}].
    \]
    \item \textbf{塔式性质}：如果 $\mathcal{H}\subseteq\mathcal{G}$，则
    \[
        \mathbb{E}[\mathbb{E}[X\mid\mathcal{G}]\mid\mathcal{H}]
        =\mathbb{E}[X\mid\mathcal{H}].
    \]
    \item \textbf{全期望公式}：
    \[
        \mathbb{E}[\mathbb{E}[X\mid\mathcal{G}]]=\mathbb{E}[X].
    \]
\end{enumerate}
\end{theorem}

\newpage

\begin{example}[对有限划分条件化]
设 $\mathcal{G}$ 由概率均为正的划分 $A_1,\dots,A_n$ 生成。则
\[
    \mathbb{E}[X\mid\mathcal{G}]
    =\sum_{k=1}^n
    \frac{\mathbb{E}[X\mathbf{1}_{A_k}]}{\mathbb{P}(A_k)}
    \mathbf{1}_{A_k}.
\]
条件期望在每个信息单元 $A_k$ 上为常数，因为 $\mathcal{G}$ 无法区分同一单元内部的不同结果。
\end{example}

\begin{theorem}[最小二乘刻画]
如果 $X\in L^2$，且 $Y=\mathbb{E}[X\mid\mathcal{G}]$，那么对每个平方可积的 $\mathcal{G}$-可测随机变量 $Z$，都有
\[
    \mathbb{E}[(X-Y)^2]
    \leq
    \mathbb{E}[(X-Z)^2].
\]
\end{theorem}

\begin{proof}
条件期望的定义性质推出
\[
    \mathbb{E}[(X-Y)(Y-Z)]=0.
\]
因此
\[
    \mathbb{E}[(X-Z)^2]
    =\mathbb{E}[(X-Y)^2]
     +\mathbb{E}[(Y-Z)^2],
\]
从而得到结论。
\end{proof}

\begin{remark}
这就是内积空间中正交投影的概率论对应物。$\sigma$-代数 $\mathcal{G}$ 决定哪些随机变量是可观测的，而条件期望把 $X$ 投影到这些信息上。
\end{remark}

\section{鞅与停时}

鞅刻画条件漂移为零的随机过程。它并不意味着过程保持常数，也不意味着样本路径稳定；它只意味着在给定当前信息后，下一时刻取值的条件期望等于当前取值。这一区别非常重要：公平博弈可以剧烈波动，尽管其中不存在可预测的获利机会。

\begin{example}[中心化部分和]
设 $X_1,X_2,\dots$ 是相互独立的可积随机变量，具有共同均值 $\mu$，定义
\[
    M_n=\sum_{k=1}^n(X_k-\mu).
\]
对自然滤过 $\mathcal{F}_n=\sigma(X_1,\dots,X_n)$，有
\[
    \mathbb{E}[M_{n+1}\mid\mathcal{F}_n]
    =M_n+\mathbb{E}[X_{n+1}-\mu]
    =M_n.
\]
因此 $(M_n)$ 是鞅。
\end{example}

\begin{example}[变换后的随机游走]
对简单对称随机游走 $S_n$，过程
\[
    M_n=S_n^2-n
\]
是鞅。事实上，
\[
    S_{n+1}^2=(S_n+\xi_{n+1})^2
    =S_n^2+2S_n\xi_{n+1}+1,
\]
对过去信息条件化得到
\[
    \mathbb{E}[S_{n+1}^2-(n+1)\mid\mathcal{F}_n]
    =S_n^2-n.
\]
\end{example}

\begin{definition}[停时]
随机时刻 $\tau:\Omega\to\{0,1,2,\dots,\infty\}$ 称为关于 $(\mathcal{F}_n)$ 的\textbf{停时}，如果对每个 $n$，
\[
    \{\tau\leq n\}\in\mathcal{F}_n.
\]
也就是说，利用时刻 $n$ 已经获得的信息，可以判断停止是否已经发生。
\end{definition}

\begin{theorem}[有界停时的可选停止定理]
设 $(M_n)$ 是鞅，$\tau$ 是被某个确定整数 $N$ 控制的停时。则
\[
    \mathbb{E}[M_\tau]=\mathbb{E}[M_0].
\]
\end{theorem}

\begin{proof}
写成
\[
    M_\tau=M_0+\sum_{k=1}^N
    (M_k-M_{k-1})\mathbf{1}_{\{\tau\geq k\}}.
\]
事件 $\{\tau\geq k\}$ 属于 $\mathcal{F}_{k-1}$，所以
\[
\mathbb{E}[(M_k-M_{k-1})\mathbf{1}_{\{\tau\geq k\}}]
=\mathbb{E}\left[
\mathbf{1}_{\{\tau\geq k\}}
\mathbb{E}[M_k-M_{k-1}\mid\mathcal{F}_{k-1}]
\right]=0.
\]
求和即得定理。
\end{proof}

\begin{note}
可选停止定理并不对任意无界停时成立。还需要有界增量、可积性或一致可积性等附加条件。“公平博弈在自行选择停止时刻后仍然公平”这一非正式说法，只有在停止规则满足适当条件时才正确。
\end{note}

\section{Markov 链}

Markov 链是一类未来只通过当前状态依赖过去的随机过程。这并不等同于独立性：相邻状态通常是相关的。Markov 性质的含义是，当前状态已经包含了过去信息中所有与预测下一状态有关的部分。

\begin{definition}[时齐 Markov 链]
取值于可数状态空间 $S$ 的离散时间过程 $(X_n)_{n\geq0}$ 称为\textbf{时齐 Markov 链}，如果在条件概率有定义时，
\[
\mathbb{P}(X_{n+1}=j\mid X_n=i,X_{n-1}=i_{n-1},\dots,X_0=i_0)
=p_{ij}.
\]
矩阵
\[
    P=(p_{ij})_{i,j\in S}
\]
称为\textbf{转移矩阵}。
\end{definition}

\begin{proposition}[Chapman--Kolmogorov 方程]
令 $p_{ij}^{(n)}=\mathbb{P}(X_n=j\mid X_0=i)$，则
\[
    p_{ij}^{(m+n)}=
    \sum_{k\in S}p_{ik}^{(m)}p_{kj}^{(n)}.
\]
矩阵形式为
\[
    P^{(m+n)}=P^{(m)}P^{(n)},
\]
因此对时齐链有 $P^{(n)}=P^n$。
\end{proposition}

\begin{definition}[平稳分布]
状态空间 $S$ 上的概率向量 $\pi$ 称为\textbf{平稳分布}，如果
\[
    \pi P=\pi.
\]
如果 $X_0$ 的分布为 $\pi$，那么每个 $X_n$ 的分布都为 $\pi$。
\end{definition}

\begin{example}[两状态链]
设
\[
P=
\begin{pmatrix}
1-a&a\\
b&1-b
\end{pmatrix},
\qquad a,b\in(0,1).
\]
联立 $\pi P=\pi$ 与 $\pi_0+\pi_1=1$，得到
\[
    \pi=\left(\frac{b}{a+b},\frac{a}{a+b}\right).
\]
某状态的流入转移率相对于流出转移率越大，它的长期平稳权重就越大。
\end{example}

\begin{theorem}[有限不可约非周期链]
有限 Markov 链若不可约且非周期，则存在唯一平稳分布 $\pi$，并且对每个初始状态 $i$，
\[
    p_{ij}^{(n)}\to\pi_j
\]
当 $n\to\infty$ 时成立。
\end{theorem}

\begin{remark}
不可约意味着每个状态最终都能到达其他任意状态；非周期则排除了被迫遵循的循环节律。二者共同保证该链逐渐遗忘初始状态并收敛到平衡。
\end{remark}

\section{Poisson 过程与 Brownian 运动}

两类连续时间过程在概率论及其应用中反复出现。Poisson 过程刻画以恒定平均速率随机发生的事件，Brownian 运动刻画连续的随机涨落。它们的样本路径外观完全不同：Poisson 路径以整数幅度跳跃，而 Brownian 路径几乎处处连续，却几乎必然处处不可微。

\subsection{Poisson 过程}

\begin{definition}[Poisson 过程]
如果过程 $(N_t)_{t\geq0}$ 满足：
\begin{enumerate}
    \item $N_0=0$；
    \item 具有独立增量；
    \item 具有平稳增量；
    \item 对 $s,t\geq0$，
    \[
        N_{t+s}-N_s\sim\operatorname{Poisson}(\lambda t),
    \]
\end{enumerate}
则称它是速率为 $\lambda>0$ 的\textbf{Poisson 过程}。
\end{definition}

\begin{proposition}
对 Poisson 过程，
\[
    \mathbb{E}[N_t]=\lambda t,
    \qquad
    \operatorname{Var}(N_t)=\lambda t.
\]
此外，等待第一次事件发生的时间 $T_1$ 服从参数为 $\lambda$ 的指数分布：
\[
    \mathbb{P}(T_1>t)=\mathbb{P}(N_t=0)=e^{-\lambda t}.
\]
\end{proposition}

\begin{example}
如果电话以每小时 $3$ 次的速率按照 Poisson 过程到达，那么两小时内恰好接到五个电话的概率为
\[
    \mathbb{P}(N_2=5)=e^{-6}\frac{6^5}{5!}.
\]
等待下一个电话超过二十分钟的概率为
\[
    e^{-3(1/3)}=e^{-1}.
\]
\end{example}

\subsection{Brownian 运动}

\begin{definition}[标准 Brownian 运动]
如果过程 $(B_t)_{t\geq0}$ 满足：
\begin{enumerate}
    \item $B_0=0$ 几乎必然成立；
    \item 具有独立增量；
    \item 对 $0\leq s<t$，有 $B_t-B_s\sim N(0,t-s)$；
    \item 样本路径几乎必然连续，
\end{enumerate}
则称它是\textbf{标准 Brownian 运动}。
\end{definition}

\begin{proposition}
对 Brownian 运动，
\[
    \mathbb{E}[B_t]=0,
    \qquad
    \operatorname{Var}(B_t)=t,
    \qquad
    \operatorname{Cov}(B_s,B_t)=\min(s,t).
\]
\end{proposition}

\begin{proof}
设 $s\leq t$。由于 $B_t=B_s+(B_t-B_s)$，且增量与 $B_s$ 独立并具有零均值，
\[
    \mathbb{E}[B_sB_t]
    =\mathbb{E}[B_s^2]
     +\mathbb{E}[B_s]\mathbb{E}[B_t-B_s]
    =s.
\]
因此协方差为 $s=\min(s,t)$。
\end{proof}

\newpage

\begin{example}[Brownian 鞅]
过程 $(B_t)$ 关于其自然滤过是鞅。另一个重要鞅是
\[
    B_t^2-t.
\]
减去 $t$ 补偿了随时间累积的方差，正如简单随机游走中 $S_n^2-n$ 的作用一样。
\end{example}

\begin{remark}
尺度性质
\[
    (B_{ct})_{t\geq0}
    \overset{d}{=}
    (\sqrt{c}B_t)_{t\geq0}
\]
表明时间相对于空间按平方尺度变化。同样的尺度关系也出现在热方程中，并解释了 Brownian 运动与扩散之间的密切联系。
\end{remark}


\section{结语}

概率论从有限样本空间中的计数问题出发，最终发展为建立在测度论基础上的严密理论。随机变量是可测函数，期望是积分，独立性是乘积分布结构，条件期望是给定信息下的最佳预测。大数定律和中心极限定理解释了随机现象中的稳定性与正态近似，而随机过程则把概率论推向时间演化和动态系统。

\textbf{关键词：} 概率空间，随机变量，期望，方差，独立性，条件期望，收敛方式，大数定律，中心极限定理，随机过程。
